\chapter{재현 실험 프로토콜}
\label{app:experiments}

\section{재현의 최소 단위}

해석 실험은 그림이 아니라 실행 조건과 raw 결과가 함께 있어야 재현된다. 모든 실험은
다음 manifest를 저장한다.

\begin{itemize}
  \item repository commit과 변경 여부;
  \item model repository, revision과 weight file digest;
  \item tokenizer repository, revision, chat template와 실제 token IDs;
  \item framework, CUDA 또는 Metal, kernel, device와 dtype;
  \item random seed, deterministic 설정과 batch 순서;
  \item dataset source, license, sample ID와 필터링 규칙;
  \item hook 이름, tensor shape와 저장 시점;
  \item intervention, baseline, metric과 aggregation code;
  \item 표에 쓰기 전 sample-level raw value와 실패 log.
\end{itemize}

artifact directory는 실험 입력과 결과를 덮어쓰지 않도록 run ID별로 분리한다. 다음 구조는
필수 규칙이 아니라 권장 예시다.

\begin{lstlisting}[language={},caption={권장 실험 artifact 구조}]
runs/<run-id>/
  manifest.json
  config.yaml
  tokenized_inputs.jsonl
  metrics.parquet
  interventions.parquet
  figures/
  logs/
  environment.txt
\end{lstlisting}

\section{Lab 1: 손으로 계산한 attention과 코드의 일치}

첫 실험의 목표는 길이 3, head dimension 2인 causal attention을 한 숫자도 건너뛰지 않고
재현하는 것이다. 정수 또는 소수 한 자리의 Q, K, V를 직접 정하고 score, scale, mask,
row-wise softmax와 weighted sum을 종이에 계산한다.

성공 기준은 독립 구현 두 개가 같은 결과를 내는 것이다. 하나는 명시적인 for-loop로,
다른 하나는 matrix multiplication으로 작성하고 double precision에서 각 원소 오차를
$10^{-10}$ 이하로 확인한다. mask를 softmax 전과 후에 잘못 적용한 negative test도 넣어
차이를 기록한다.

\section{Lab 2: 최소 decoder-only Transformer 구현}

두 번째 실험의 목표는 high-level Transformer module을 호출하지 않고 한 block의 모든
tensor를 직접 만드는 것이다. 최소 구현에는 token embedding, RMSNorm, RoPE, causal
MHA 또는 GQA, SwiGLU, residual addition, final norm과 LM head가 포함되어야 한다.

성공 기준은 모든 중간 shape와 causal invariant가 assertion을 통과하는 것이다. 미래 token
ID만 바꾼 두 sequence를 넣었을 때 그 이전 position의 logits가 허용 오차 안에서 같아야
한다. 각 attention row의 합, masked weight, RMS denominator와 final vocabulary dimension도
검사한다.

hook 출력은 parameter와 activation을 분리하여 두 종류의 값을 혼동하지 않게 해야 한다.
다음 의사코드는 최소 shape trace를 보여 준다.

\begin{lstlisting}[language=Python,caption={한 block의 shape 추적 의사코드}]
x = embed(token_ids)                 # [B, T, d]
r = rmsnorm_attn(x)                  # [B, T, d]
q, k, v = project_and_split(r)       # [B, Hq/Hkv, T, dh]
q, k = apply_rope(q, k, positions)
a = causal_softmax(q @ k.transpose(-1, -2) / sqrt(dh))
x = x + merge_heads(a @ repeat_kv(v)) @ W_o
r = rmsnorm_mlp(x)
x = x + (silu(r @ W_gate) * (r @ W_up)) @ W_down
logits = final_norm(x) @ W_unembed.T # [B, T, V]
\end{lstlisting}

\section{Lab 3: 공개 reference implementation과 대조}

세 번째 실험의 목표는 직접 구현한 block을 실제 공개 checkpoint의 reference code와
대조하는 것이다. 같은 weight를 load하고 tokenizer output, embedding, 각 normalization,
Q/K/V, RoPE 뒤 tensor, attention output, MLP output, residual과 logits를 hook한다.

성공 기준은 dtype별 허용 오차와 첫 불일치 지점을 자동 보고하는 것이다. 마지막 logit만
비슷하다고 통과시키지 않으며, layer별 최대 절대오차와 상대오차를 저장한다. fused kernel로
중간값을 직접 꺼낼 수 없으면 unfused reference path와 비교한다.

불일치의 흔한 원인은 transposed weight, RoPE의 even/odd packing, GQA의 head mapping,
RMSNorm $\epsilon$, bias, position offset와 tokenizer special token이다. 원인을 고친 뒤에는
해당 오류를 재현하는 regression test를 남긴다.

\section{Lab 4: prefill과 KV-cache decode의 동등성}

네 번째 실험의 목표는 full recomputation과 cached decode가 같은 next-token distribution을
계산하는지 검증하는 것이다. prompt의 마지막 logit을 full pass로 얻고, token을 하나씩
넣는 cached path의 각 step logit과 비교한다.

성공 기준은 동일 dtype과 deterministic kernel에서 정한 오차 안에 있고 greedy token이
모든 step에서 일치하는 것이다. cache의 layer, KV head, sequence와 head-dimension 축을
기록하고, RoPE position이 cache length만큼 offset되는지 검사한다.

\section{Lab 5: probe는 정보 존재만 보고한다}

다섯 번째 실험의 목표는 representation에서 label을 복원하는 증거와 모델 사용의 증거를
분리하는 것이다. factual prompt의 last-subject와 last-position residual을 layer별로 모아
relation 또는 정답 class에 대한 regularized linear probe를 학습한다.

성공 기준은 held-out entity split, random-label control과 majority baseline을 포함하는 것이다.
같은 entity가 train과 test의 paraphrase에 동시에 들어가지 않도록 split 단위를 entity로
잡는다. probe accuracy, calibration, selectivity와 confidence interval을 함께 보고한다.

후속 intervention은 probe direction이 원 모델 행동에 미치는 효과를 따로 시험한다. direction을
더하고 빼는 scale sweep에서 target logit, KL divergence, unrelated-task degradation과 norm
변화를 기록한다. 효과가 있어도 자연 입력에서 원 모델이 같은 방향을 사용한다는 결론은
path-level 검증 전까지 보류한다.

\section{Lab 6: induction circuit 재현}

여섯 번째 실험의 목표는 반복 token pattern에서 induction score와 인과 효과를 함께
측정하는 것이다. 무작위 token 열의 앞부분을 뒤에서 반복하고, 현재 token과 같은 token의
바로 다음 position에 attention하는 pattern을 head별로 계산한다.

성공 기준은 pattern score가 높은 head의 ablation이 반복 구간의 next-token loss를 선택적으로
악화시키는지 확인하는 것이다. 비반복 control 구간, 크기 일치 random head와 previous-token
head 후보를 함께 비교한다. 큰 attention pattern만으로 induction head label을 확정하지
않는다.

\section{Lab 7: activation patching의 방법 민감도}

일곱 번째 실험의 목표는 clean/corrupted pair에서 localization map이 metric과 corruption에
얼마나 의존하는지 측정하는 것이다. subject 또는 relation 하나만 바꾼 factual prompt pair를
만들고 residual pre, attention output과 MLP output을 position-layer별로 patch한다.

성공 기준은 최소 세 metric과 세 baseline의 순위 상관을 보고하는 것이다. target logit,
logit difference와 probability를 비교하고, zero, mean과 resample replacement를 비교한다.
sample별 recovery가 분모가 작은 경우를 별도 표시하고, 평균 heatmap만 보고하지 않는다.

\section{Lab 8: 사실 회상의 세 단계 재현}

여덟 번째 실험의 목표는 Geva 등의 세 단계를 같은 모델에서 재현한 뒤 architecture transfer를
시험하는 것이다. 첫 단계에서는 last-subject projection의 attribute rate와 early MLP
cancellation을, 둘째 단계에서는 relation-to-last attention knockout을, 셋째 단계에서는
subject-to-last upper attention과 attribute logit 기여를 측정한다.

성공 기준은 원 논문의 GPT-2 XL 또는 GPT-J 조건을 먼저 재현하고, Llama/Qwen 계열의 결과를
별도 실험으로 보고하는 것이다. layer 번호는 절대값과 전체 깊이에 대한 비율을 함께 기록하며,
다른 architecture에서 실패해도 원 결과의 재현 실패와 generalization 실패를 구분한다.

\section{Lab 9: 2510.09033의 분류와 검출 결과 감사}

아홉 번째 실험의 목표는 FA/AH/UH label이 subject-reliance threshold에 얼마나 민감한지
검증하는 것이다. greedy BF16 생성, factual correctness 판정, subject-to-last attention
blocking, 원 분포와 개입 분포의 JS divergence, FA 평균 threshold의 순서로 원 절차를
재현한다 \citep{cheang2026knowledge}.

성공 기준은 threshold sweep과 독립 label audit를 포함하는 것이다. FA 평균뿐 아니라
quantile과 held-out calibration threshold를 비교하고, AH/UH 표본 이동률을 보고한다.
probe train/test split은 entity와 template leakage를 막으며, correctness judge의 일부
표본은 사람이 blind review한다.

원 연구의 conclusion과 label construction의 결합을 따로 검사한다. AH/UH가 subject-flow
intervention으로 정의되므로 같은 flow metric에서 나타나는 차이는 부분적으로 예상된다.
새로운 검증에는 label에 직접 쓰지 않은 intervention, paraphrase robustness와 외부 factual
verification을 추가한다.

\section{Lab 10: ground-truth circuit에서 방법 평가}

열 번째 실험의 목표는 자연 모델에 적용하기 전에 해석 도구의 탐지 오류를 알려진 program에서
측정하는 것이다. Tracr로 token frequency, sorting 또는 parenthesis checking program을
컴파일하고 source code를 보지 않는 분석 단계와 ground-truth 비교 단계를 분리한다.

성공 기준은 component와 edge recovery의 precision, recall, faithfulness와 description length를
함께 보고하는 것이다. probe, activation patching, ACDC, attribution patching과 SAE를 같은
node granularity에서 비교한다. 알고리즘 이름을 맞히는 것뿐 아니라 unseen input과
intervention output을 예측하는지도 평가한다.

\section{실험 중단과 실패 보고 규칙}

실패한 재현은 조건을 충분히 기록하면 유효한 결과다. model revision, token mismatch,
memory 부족, numerical divergence, 불안정한 metric과 논문에 없는 구현 세부사항을 실패
유형으로 나누어 남긴다.

결론을 바꾸는 분석 선택은 결과를 본 뒤 조용히 수정하지 않는다. threshold, layer window,
metric과 sample filter를 사전에 config에 저장하고, 사후 탐색은 exploratory analysis로
표시한다. 여러 선택을 시험했다면 성공한 설정만 골라 보고하지 않고 전체 sensitivity를
제시한다.

\begin{studybox}[공통 통과 기준]
실험은 다른 사람이 새 환경에서 manifest만으로 실행할 수 있고, raw sample-level 결과에서
책의 표와 그림을 다시 만들 수 있으며, negative control과 적어도 하나의 합리적인 대안
설정에서 핵심 결론이 유지될 때 통과한다.
\end{studybox}
